\documentclass[lettersize,journal]{IEEEtran}
\usepackage{amsmath,amsfonts}
\usepackage{algorithmic}
\usepackage{algorithm}
\usepackage{array}
\usepackage[caption=false,font=normalsize,labelfont=sf,textfont=sf]{subfig}
\usepackage{textcomp}
\usepackage{stfloats}
\usepackage{url}
\usepackage{verbatim}
\usepackage{graphicx}
\usepackage{cite}
\hyphenation{op-tical net-works semi-conduc-tor IEEE-Xplore}

\begin{document}

\title{Optimizing Large-Scale General Matrix Multiplication on Heterogeneous DSPs via Hierarchical Pipelining and Scalable Communication}

\author{Antigravity GEMM Group, National Supercomputing Center
\thanks{Manuscript received June 15, 2026. This work was supported by the key research and development program of the MT-3000 high-performance computing system validation group.}}

% The paper headers
\markboth{Journal of High Performance Computing Class Files,~Vol.~15, No.~6, June~2026}%
{Antigravity GEMM Group \MakeLowercase{\textit{et al.}}: Optimizing Large-Scale General Matrix Multiplication on Heterogeneous Multi-Core DSPs via Hierarchical Pipelining and Scalable Communication}

\IEEEpubid{0000--0000/00\$00.00~\copyright~2026 IEEE}

\maketitle

\begin{abstract}
High-performance General Matrix Multiplication (GEMM) is a primary computational engine for exascale scientific applications. Modern many-core architectures, such as the domestic MT-3000 processor, present exceptional computational density by utilizing software-managed scratchpads and VLIW-SIMD hardware. However, deploying GEMM at a massive scale on heterogeneous systems faces severe challenges: network contention during inter-process data exchange leads to communication bottlenecks, while manually tuning parameters across explicit memories yields suboptimal hardware utilization. This paper proposes a highly scalable, optimized complex-valued GEMM engine on the MT-3000. Under a hierarchical 2D data-partitioning model, we introduce a scalable group-broadcast communication scheme with asymmetric workload mapping, dynamically balancing networking overhead between CPU group leaders and memory-bound accelerators. At the device core level, we introduce a tri-level pipelining framework combining A/B double buffering and C triple buffering in array memories, which hides explicit Direct Memory Access (DMA) latencies. VLIW instructions are optimized via 16-way SIMD assembly-like vector operations to exploit outer-product arithmetic structures. Finally, a constraint-satisfying surrogate auto-tuner is developed by using offline Random Forest regressors to navigate the multi-dimensional blocking parameters. Experimental results on the MT-3000 supercomputer show that our framework maintains a strong-scaling efficiency of $82.1\%$ at 256 processes, achieving a sustained throughput of 21.6 TFLOPS (over $91.5\%$ of the theoretical peak), which represents a $5.14\times$ speedup over standard flat-broadcast SUMMA implementations.
\end{abstract}

\begin{IEEEkeywords}
Heterogeneous optimization, Many-core DSP, Complex General Matrix Multiplication, Explicit memory scheduling, Auto-tuning, Random Forest surrogate.
\end{IEEEkeywords}

\section{Introduction}

\IEEEPARstart{G}{eneral} Matrix Multiplication (GEMM) is the fundamental subprogram in high-performance computing (HPC), serving as the core computational engine for a vast spectrum of scientific and engineering workloads. From traditional numerical simulations, such as Computational Fluid Dynamics (CFD) and the Method of Moments (MoM) in electromagnetics, to modern deep learning frameworks, the overall performance of these applications is essentially dictated by the efficiency of the underlying GEMM kernels. Because GEMM frequently accounts for the vast majority of total execution cycles in these compute-intensive tasks, even marginal improvements in its throughput can lead to significant gains in application-level performance. As HPC systems transition to the exascale era, heterogeneous architectures based on Digital Signal Processors (DSPs), such as the MT-3000, have emerged as a promising solution. These processors are specifically designed to deliver massive floating-point throughput and superior energy efficiency for matrix-intensive operations; however, their unique architectural disparities from conventional CPUs and GPUs pose significant challenges in unleashing their full potential for GEMM.

The MT-3000 processor exemplifies the paradigm of extreme many-core heterogeneity designed for exascale-level computation. It features a hierarchical multi-zone architecture that integrates general-purpose processing zones for system management with massive acceleration clusters for high-density arithmetic~\cite{ref17}. Recent work has further demonstrated the effectiveness of this architecture in large-scale LLM training scenarios~\cite{ref19}. The most distinctive characteristic of this architecture is its shift from a conventional hardware-managed cache hierarchy to a purely software-managed scratchpad memory system, coupled with a VLIW-SIMD execution model. By eschewing traditional caches in favor of explicit on-chip memories and high-bandwidth interconnects, the MT-3000 achieves exceptional arithmetic density and energy efficiency. 

Despite its immense potential, optimizing GEMM on the MT-3000 for large-scale environments spanning hundreds or thousands of processes faces two fundamental challenges.
First, achieving high scalability is hindered by the "communication storm" bottleneck in large-scale multi-process environments. As the system scales, standard flat matrix distribution patterns trigger severe network contention, leading to compute core idleness and a sharp decline in strong-scaling efficiency. To address this, we propose a Scalable Hierarchical Communication strategy based on a Leader-Member protocol. By localizing data movement and implementing an asymmetric task allocation, we enable group leaders to overlap their data-forwarding responsibilities with early-start computation, effectively transforming communication wait-time into productive compute cycles.
Second, explicit memory management and massive search space create a daunting hurdle for holistic resource orchestration. The absence of hardware-managed caches necessitates manual data movement across four hierarchical tiers, which often outpaces computation, while rigid architectural constraints and astronomical search spaces preclude manual attainment of optimal efficiency. To overcome these multifaceted issues, we develop an integrated optimization suite that orchestrates a hierarchical tri-level pipelining framework with a specialized vectorized micro-kernel, all managed by a constraint-satisfying auto-tuning model. This synergy ensures maximum compute unit saturation and deterministic load balance within the DSP clusters, effectively shielding the computation from the latencies of the explicit memory hierarchy.

The primary contributions of this paper are summarized as follows:

\begin{itemize}
\item[$\bullet$] We propose a 2D data partitioning scheme along with an associated hierarchical communication strategy. This approach ensures global load balance across large-scale process groups via asymmetric workload distribution.
\item[$\bullet$] We develop a tri-level hierarchical pipelining framework integrated with a vectorized micro-kernel design. It effectively hides the latencies of explicit multi-level memory structures by overlapping data movement with computation.
\item[$\bullet$] We introduce an auto-tuning model for GEMM parameters based on surrogate modeling. It enables rapid identification of optimal configurations while strictly adhering to hardware alignment and capacity constraints.
\end{itemize}

\section{BACKGROUND}

\subsection{MT-3000 Many-core Architecture}
    
The MT-3000 is a heterogeneous many-core processor designed for exascale computing through a multi-zone architecture. It partitions resources into a general-purpose zone containing sixteen CPU cores and four independent acceleration clusters. While the CPU cores manage system orchestration and inter-process MPI communication, the acceleration clusters function as high-throughput engines for compute-intensive tasks. Each cluster contains 24 acceleration arrays supported by on-chip Global Shared Memory (GSM) and High-Bandwidth Shared Memory (HBSM). This clear separation allows the general-purpose zone to coordinate global data movement while the acceleration clusters focus on localized, high-density matrix computations.


The fundamental unit of parallelism within each cluster is the acceleration array, which utilizes a specialized control-compute micro-architecture. Each array integrates a single Scalar Processing Unit (SPU) and sixteen Vector Processing Units (VPUs) working in lock-step. The SPU governs instruction flow and scalar address calculations, while the sixteen VPUs serve as the primary computational engines using a Very Long Instruction Word (VLIW) design. This micro-architecture enables high-density Single Instruction Multiple Data (SIMD) execution, capable of issuing multiple vector instructions—such as multiply-add and load/store—in a single cycle. A crucial feature is the Scalar-Vector Register (SVR), which allows the SPU to broadcast scalar elements directly to all VPUs. This provides the hardware-level foundation for outer-product-based matrix multiplication, where a scalar from one matrix can be efficiently broadcast to a vector segment of another.


The memory system of MT-3000 follows a multi-tier, software-managed hierarchy consisting of off-chip DDR memory, cluster-shared GSM, and core-private Array Memory (AM) and Scalar Memory (SM). In contrast to processors with hardware-coherent cache hierarchies, the MT-3000 requires all data transfers between memory tiers to be explicitly orchestrated through Direct Memory Access (DMA) engines. This tiered structure and explicit memory management provide the underlying infrastructure for precise data staging. By manually coordinating the movement of blocks across the DDR, GSM, and AM, GEMM-like operators can be optimized to achieve high hardware utilization and sustained performance.

\subsection{Algorithms for Large-scale Parallel GEMM}

In the field of large-scale parallel computing, the performance of General Matrix Multiplication (GEMM) is fundamentally determined by the strategies used for data partitioning and distribution, which dictate the ratio of communication to computation. Parallel GEMM algorithms are traditionally categorized into one-dimensional (1D), two-dimensional (2D), and three-dimensional (3D) decomposition. 1D partitioning segments matrices into row or column stripes; while straightforward, it lacks scalability in many-core systems due to limited concurrency and excessive communication volumes. Although 3D decomposition offers higher theoretical parallelism by distributing data across a three-dimensional process grid, it often incurs significant memory overhead and topological complexity. Consequently, 2D decomposition has emerged as the most widely adopted paradigm for large-scale systems, striking an optimal balance between memory efficiency and scalability.

Under the 2D decomposition framework, the process space is logically organized into a grid. The output matrix $C$ is partitioned into sub-blocks, with each process assigned to compute a specific local block $C_{i,j}$. The computation of this sub-block is inherently tied to the acquisition of the corresponding row-blocks from matrix A and column-blocks from matrix B. This spatial partitioning effectively decouples the global computational workload into independent localized tasks that can be executed concurrently across a large number of processes. This 2D partitioning logic establishes the essential relationship between data placement, communication overhead, and computational concurrency, providing the structural basis for matrix operations on many-core architectures.

As a representative implementation of the 2D partitioning logic, the Scalable Universal Matrix Multiplication Algorithm (SUMMA)~\cite{ref3} has become a cornerstone of modern parallel linear algebra libraries. The core principle of SUMMA is a broadcast-driven communication model, often referred to as "broadcast-and-multiply." It operates by asynchronously broadcasting segments of matrix A and matrix B along the rows and columns of the process grid, allowing compute units to progressively accumulate local products. A key advantage of SUMMA is its adaptability to various process grid shapes and its ability to hide communication latency through pipelined operations. This broadcast-driven distribution logic, built on 2D partitioning, provides a rigorous theoretical foundation for parallel GEMM. It serves as a baseline for further advances on heterogeneous processors, such as the implementation of hierarchical communication and scalable load-balancing strategies.

\section{RELATED WORK}
\subsection{GEMM Optimization on Heterogeneous DSP Architectures}

Over the past few decades, the community has established the Goto algorithm~\cite{ref1} as the gold standard for General Matrix Multiply (GEMM) optimization on conventional cache-coherent CPUs. This paradigm, utilized by mainstream libraries such as OpenBLAS and Intel MKL~\cite{ref18}, relies on hardware caches to hide memory latency through multi-level tiling and data packing. However, the architecture of heterogeneous many-core processors and digital signal processors (DSPs), such as the MT-3000, represents a fundamental departure from this model. On these systems, transparent cache hierarchies are replaced by software-managed scratchpad memories (SPM), requiring a shift from implicit cache-centric optimization to explicit Direct Memory Access (DMA) orchestration. Early research on the Cell Broadband Engine—a pioneering heterogeneous architecture with master-slave processing elements—demonstrated that near-peak performance could be achieved by manually overlapping DMA transfers with SIMD execution~\cite{ref2}. This explicit-control methodology was later extended to massive-scale many-core systems like the Sunway SW26010, where researchers utilized 
Sunway SW26010, where researchers utilized register-level tiling and slave-core local memory management to overcome the absence of hardware caches~\cite{ref9}. More recently, the Sunway Math Library has further advanced performance through superblock-based optimization techniques~\cite{ref16}. Parallel efforts on the SW26010 have also demonstrated runtime-adaptive GEMM scheduling~\cite{ref10}, highlighting the growing maturity of many-core DSP ecosystems.

The transition to multi-core DSP architectures has further specialized these techniques. For instance, studies on the TI KeyStone DSP series adapted the BLIS framework to handle hierarchical scratchpad memories, proving that heterogeneous DSPs could effectively match the performance of general-purpose processors in high-performance computing (HPC) workloads~\cite{ref4}. More recently, research specifically targeting the MT-3000 architecture has achieved significant breakthroughs in assembly-level micro-kernel efficiency by exploiting Very Long Instruction Word (VLIW) scheduling and Scalar-Vector Register (SVR) broadcasting~\cite{ref8}. While these efforts have matured the core-level execution, they often exhibit "parameter rigidity," where static analytical models fail to account for non-linear performance fluctuations such as bank contention. 

\subsection{Machine Learning-based Performance Modeling}

To overcome the overwhelming complexity of manual parameter selection, the HPC community has evolved from empirical search methods to automated performance modeling. Early efforts like ATLAS~\cite{ref5} pioneered the empirical search paradigm, successfully achieving near-peak performance on various CPUs by automating the selection of optimal blocking and unrolling factors. However, as architectures became more heterogeneous, the search space grew exponentially, rendering exhaustive search-based methods computationally prohibitive. Modern deep learning compilers like Ansor~\cite{ref6} addressed this by introducing machine learning-based cost models to predict program performance, demonstrating that gradient-boosted trees can outperform manual expert tuning by up to 2× on mainstream CPUs and GPUs. Despite these achievements, these frameworks often struggle with the non-linear performance landscapes of DSPs, where instruction-level VLIW scheduling and hardware-specific bank conflicts create unique optimization barriers.

Recent research has pivoted toward surrogate modeling to navigate high-dimensional parameter spaces with minimal experimental overhead. Studies have proven that Random Forest models are particularly effective in capturing intricate hardware-software interactions from sparse experimental data, providing a robust and data-driven alternative to rigid analytical models~\cite{ref7}. Building on this, Qi et al.~\cite{ref8} recently achieved a milestone in auto-tuning efficiency for the MT-3000 by demonstrating a constraint-driven framework for GEMM-like operators at SC 2025, leveraging hardware-guided pruning to reduce the search space. While this approach significantly accelerated kernel-level optimization, most existing auto-tuners still focus exclusively on single-node execution. Our work extends these successes by utilizing a Random Forest-based surrogate model to jointly optimize local kernel efficiency and global process-level load balancing, providing a scalable solution for massive-scale heterogeneous GEMM deployments.
 

\section{DATA ORCHESTRATION AND COMMUNICATION }
\subsection{Hierarchical Data Partitioning}
To maximize computational throughput and eliminate memory bandwidth bottlenecks for complex-valued General Matrix Multiplication (GEMM) $C = A \times B$, where $A \in \mathbb{C}^{M \times K}$, $B \in \mathbb{C}^{K \times N}$, and $C \in \mathbb{C}^{M \times N}$, we design a hierarchical, two-level 2D data partitioning scheme. This scheme encompasses a macro-level inter-process partition across the host CPU nodes, and a micro-level intra-device scratchpad partitioning within the many-core DSP acceleration arrays.

At the host level, the global process grid consists of $P$ MPI processes, logically arranged to facilitate column-oriented partitioning for the matrix $B$ and row-oriented partitioning for the matrix $A$. The processes are clustered into $G$ subgroups, with each subgroup consisting of $P_g = P / G$ processes. Matrix $B$ is evenly partitioned along its columns into $G$ vertical blocks (stripes). The $g$-th process subgroup is assigned the task of processing the column range $[start\_col[g], start\_col[g+1])$, representing a local column width of $N_{local} \approx N / G$. Meanwhile, matrix $A$ is partitioned along the row dimension. To alleviate the communication and coordination overhead incurred by the group leader (which handles inter-process data routing), we introduce an asymmetric row allocation mechanism:
The group leader of each subgroup (having a local group rank of $0$) is assigned $M_{leader} = groupm$ rows of matrix $A$. To manage the remainder rows when the remaining matrix size is not perfectly divisible, we formulate the row allocation for any process $t$ within group $g$ (where local group rank is $t_g = t \pmod{P_g}$) as:
\begin{equation}
\text{local\_M}[t] = \begin{cases} 
M_{leader}, & t_g = 0 \\
M_{member} + \delta_t, & t_g > 0
\end{cases}
\end{equation}
where $M_{member} = \lfloor (M - M_{leader}) / (P_g - 1) \rfloor$, and the remainder is $M_{remaining\_rows} = (M - M_{leader}) \pmod{P_g - 1}$. The active offset indicator $\delta_t$ is defined as:
\begin{equation}
\delta_t = \begin{cases}
1, & t_g \le M_{remaining\_rows} \\
0, & t_g > M_{remaining\_rows}
\end{cases}
\end{equation}
This mathematically guarantees deterministic global row load balance across all MPI nodes while dynamically keeping the networking leaders slightly under-burdened by calculations.

At the device level (intra-DSP), the memory space is software-managed through explicit Direct Memory Access (DMA) transactions. The local matrix fragments allocated to each MPI process must be tiled to fit inside the extremely constrained on-chip scratchpad memories. This requires three distinct tiling parameters: $M_{block}$, $K_{block}$, and $N_{block}$.
For any process, its local row-block of matrix $A$ is partitioned into tiles of size $M_{block} \times K_{block}$. These tiles are dynamically loaded into the $64\text{ KB}$ Scalar Memory (SM), which is shared among all Vector Processing Units (VPUs) and governed by the SPU. Conversely, the column-blocks of matrix $B$ are tiled into fragments of size $K_{block} \times N_{block}$ and loaded into the $768\text{ KB}$ Array Memory (AM), which is private to each VPU. The intermediate matrices of $C$ are also allocated as $M_{block} \times N_{block}$ blocks within the AM. Because the AM is core-private, caching matrix $B$ and matrix $C$ in the AM allows spatial alignment with VPU lanes, whereas locating matrix $A$ in the SPU-controlled SM allows single-cycle broadcast of scalar multipliers. To satisfy SIMD vector alignment requirements on the MT-3000, we enforce the hardware constraint:
\begin{equation}
N_{block} \equiv 0 \pmod{16}
\end{equation}
which matches the 16-way double-precision vector register width of the execution hardware.

\subsection{Scalable Group Communication}
Standard parallel GEMM implementations, such as flat SUMMA, often suffer from severe network congestion at scale, commonly referred to as the "communication storm" bottleneck. Since all MPI processes simultaneously exchange row and column segments, the underlying network interfaces suffer from heavy collision issues, causing compute cores to idle while waiting for data. To resolve these bottlenecks, we propose a Scalable Group Communication mechanism based on a hierarchical Leader-Member protocol, which localizes the distribution traffic and overlaps communication with disk I/O.

The communication pattern is divided into two decoupled phases:
\begin{enumerate}
\item \textbf{Phase 1 (Inter-Group Asynchronous Distribution):} The root process (Rank 0) holds the complete matrix $B$ and distributes it to the designated Group Leaders. Specifically, Rank 0 issues asynchronous, non-blocking sends (\texttt{MPI\_Isend}) to transmit the corresponding column-block of matrix $B$ to each Group Leader (rank $g \times P_g$ for $g = 0, \dots, G-1$). Simultaneously, while the network interface card (NIC) handles this data transfer in the background, all processes (including regular member processes and group leaders) concurrently load their respective portions of matrix $A$ from the local filesystem or disk via asynchronous parallel I/O (\texttt{pread\_matrix}). By executing filesystem reads concurrently with inter-process network transfers, we achieve optimal overlap between host-side computation/storage operations and network communications.
\item \textbf{Phase 2 (Intra-Group Asynchronous Broadcast):} After Phase 1 completes, each subgroup creates a dedicated intra-group communicator via \texttt{MPI\_Comm\_split}, partitioning the global communicator \texttt{MPI\_COMM\_WORLD} by $group\_id$. This provides an isolated communication domain for localized collective operations. Once a Group Leader successfully completes the reception of its assigned column-strip of matrix $B$ via \texttt{MPI\_Waitall}, it acts as the subgroup root and initiates a non-blocking collective broadcast (\texttt{MPI\_Ibcast}) on the group communicator to all processes within its logical group. Crucially, because the Leader already possesses the required segments of matrix $B$ and its local chunks of matrix $A$, it does not block waiting for the group members to complete their broadcasts. Instead, the Leader proceeds directly to spawn the DSP compute threads via hyper-threading device APIs (\texttt{gemm\_dsp\_AM\_SM\_pipeline}) and begins calculation for its assigned $M_{leader}$ rows. The subgroup member processes receive the broadcasted matrix $B$ data asynchronously via \texttt{MPI\_Waitall} and then trigger their computing phases. High-level communication latency is effectively hidden, as the compute kernels of leaders and members are staggered, absorbing the transfer time and eliminating network-induced synchronization barriers.
\end{enumerate}

\section{ALGORITHM DESIGN AND OPTIMIZATION}
\subsection{Multi-tiered Pipelining}
The absolute absence of hardware-managed memory coherence or automatic caches on the MT-3000 places the responsibility of data scheduling entirely on the software layers. High computational throughput can only be achieved when data movements between private on-chip scratchpad memories and system DDR are fully overlapped with CPU/DSP arithmetic execution. To meet this challenge, we present a hierarchical tri-level pipelining framework consisting of distinct outer, middle, and inner tiers, which isolates data routing from computation.

\begin{enumerate}
\item \textbf{Outer-tier Pipeline (Inter-Process Overlapping):}
At the cluster level, the outermost stage handles the asymmetric computation and communication overlap between process groups. The Group Leader initiates asynchronous requests (\texttt{MPI\_Isend} and \texttt{MPI\_Irecv}) for column-slices of matrix $B$, while simultaneously triggering local filesystem reads (\texttt{pread\_matrix}) in parallel for matrix $A$. The communication latency of matrix $B$ distribution is effectively absorbed by the processing time of matrix $A$ storage operations, ensuring that the host CPU begins compute kernels without waiting for global network synchronization.

Specifically, the timeline is structured as follows: (1) Rank 0 posts $G-1$ non-blocking \texttt{MPI\_Isend} operations for $B$ column-stripes to peer group leaders; (2) \emph{simultaneously} all $P$ processes issue \texttt{pread\_matrix} calls to read matrix $A$ stripes from local NVMe or parallel filesystem storage; (3) these I/O operations occupy the CPU memory bandwidth for $t_{IO}$ seconds, during which the network interface card (NIC) transfers $B$ blocks over the high-speed interconnect in the background; (4) after both the I/O and \texttt{MPI\_Waitall} operations complete, the accelerated computing kernels begin execution. Under optimal scaling conditions, $t_{IO} \geq t_{comm}$, meaning the communication cost is completely hidden behind the compute-process I/O pipeline. This hierarchical communication pattern leverages topology-aware MPI grouping~\cite{ref20} to further reduce network contention.

\item \textbf{Middle-tier Pipeline (C-matrix Triple Buffering):}
Within each DSP core, write-backs of computed results from the Array Memory (AM) to off-chip DDR memory represent a major throughput hazard. Since DMA write transactions share the same on-chip data bus with read traffic, a simple direct write-back halts the matrix processing pipeline. We mitigate this through a C-matrix triple buffering strategy. Three virtual buffers (\texttt{Cx\_a1/Cy\_a1}, \texttt{Cx\_a2/Cy\_a2}, and \texttt{Cx\_a3/Cy\_a3}, representing real and imaginary components) are allocated in the private AM. These buffers operate in a rotating role-assignment scheme as the outer loop iterates over $N$-dimension column blocks. At any loop iteration, the three buffers are assigned distinct roles:
\begin{itemize}
\item \textbf{Store (retiring the oldest):} Buffer $a3$ stores its fully accumulated $C$ sub-block back to system DDR via \texttt{dma\_p2p} initiated by \texttt{chc5/chc6}.
\item \textbf{Load (prefetching the newest):} Buffer $a2$ fetches the initial zero-initialized $C$ sub-block of the next column range from DDR via \texttt{dma\_p2p} initiated by \texttt{chc3/chc4}.
\item \textbf{Compute (active accumulation):} Buffer $a1$ is actively engaged in VPU outer-product accumulation for the current $K$-tile iteration, with VPUs reading/writing via \texttt{vec\_ld} and \texttt{vec\_st}.
\end{itemize}
After each iteration boundary, the roles rotate cyclically: the former active compute buffer becomes the next store buffer, the loaded buffer becomes the new compute buffer, and the stored buffer becomes available for loading again. This circular rotation, governed by the loop variable $j$ and coordinated via non-blocking DMA channels, ensures that write-back and prefetch latencies are fully overlapped with arithmetic execution. The buffers are synchronized via hardware wait instructions (\texttt{dma\_wait}) only when pipeline transitions occur at iteration boundaries in the inner loop.

\item \textbf{Inner-tier Pipeline (A/B-matrix Double Buffering):}
At the innermost loop nesting, where the matrix multiply-accumulate micro-kernels iterate along the reduction axis $K$, latency is hidden via double buffering of matrices $A$ and $B$. We allocate two sets of registers/buffers for A in the SM (\texttt{Ax\_s1/Ay\_s1} and \texttt{Ax\_s2/Ay\_s2}) and two sets for B in the AM (\texttt{Bx\_a1/By\_a1} and \texttt{Bx\_a2/By\_a2}). While SPU/VPU arithmetic is executing on block $q$ using buffer 1, the hardware DMA engine pre-fetches block $q+1$ into buffer 2. A synchronization block enforces wait states only before the SPU switches the active register pointer. This overlapping strategy ensures that the compute pipelines are saturated with operands at every clock cycle.
\end{enumerate}

\subsection{Vectorized Micro-kernel}
The execution units of the MT-3000 rely on a 16-way SIMD Vector Processing Unit (VPU) array, which acts as the computational bedrock for the GEMM kernel. Because complex-valued general matrix multiplication involves separate computations for real and imaginary components:
\begin{equation}
C_{real} + i C_{imag} = (A_{real} \times B_{real} - A_{imag} \times B_{imag}) + i (A_{real} \times B_{imag} + A_{imag} \times B_{real})
\end{equation}
a straightforward scalar implementation suffers from severe vector registers dependency and pipeline stalls. We implement a vectorized complex-GEMM micro-kernel that maps these arithmetic operations into hardware-level SIMD intrinsics.

The core loops are designed around an outer-product computation paradigm. A column-vector of matrix $B$ is loaded into AM registers, and single elements of matrix $A$ are broadcast from SM registers to all 16 VPUs using SVR broadcast instructions (\texttt{vec\_svbcast}). Real and imaginary products are accumulated using vector intrinsics mapped exactly to VPU pipeline slots. Let $A_x, A_y$ represent vector registers loading SVR-broadcasted real and imaginary blocks of $A$, and $B_x, B_y$ represent vector registers holding $B$ column blocks in AM. Let $C_x, C_y$ represent target accumulator registers in AM. The micro-kernel coordinates complex multiply-accumulate (MAC) using the local intrinsics:
\begin{align}
\text{Real Channel (Re):}& \nonumber \\
T_{yy} &\leftarrow \texttt{vec\_muli}(A_y, B_y) \\
T_{s\_real} &\leftarrow \texttt{vec\_mulb}(A_x, B_x, T_{yy}) \\
C_x &\leftarrow \texttt{vec\_mula}(C_x, Cons, T_{s\_real}) \\
\text{Imaginary Channel (Im):}& \nonumber \\
T_{yx} &\leftarrow \texttt{vec\_muli}(A_y, B_x) \\
T_{s\_imag} &\leftarrow \texttt{vec\_mula}(A_x, B_y, T_{yx}) \\
C_y &\leftarrow \texttt{vec\_mula}(C_y, Cons, T_{s\_imag})
\end{align}
where $Cons = 1.0$ is a double-precision unity coefficient vector register; $A_x, A_y, B_x, B_y, C_x, C_y$ denote the vector register contents for the real and imaginary components mapped to the current $M_{block} \times N_{block}$ tile respectively. By executing multiplication using \texttt{vec\_muli} as an intermediate step, the subtraction $A_x B_x - A_y B_y$ is optimized as a target pipeline operation \texttt{vec\_mulb}, which executes a fused multiply-subtract operation. This sequence prevents pipeline stalls and maintains a vector register usage of 32 registers with $4\times$ loop unrolling on the VPU and SPU registers. This scheduling aligns with the 16-way double-precision vector width, reaching near-peak floating-point efficiency.

\subsection{Learning-based Auto-tuning}
Given the complexity of the hierarchical memory layout, the optimal parameters for tile dimensions ($M_{block}, K_{block}, N_{block}$), process layouts ($G, P_g$), and workload thresholds ($groupm$) form an extremely large and non-linear search space. Static compiler heuristics often fail due to hard-to-model hardware interactions (e.g., bank conflicts in AM or bus congestion). We resolve this orchestration problem via a constraint-guided, learning-based auto-tuning framework.

\begin{enumerate}
\item \textbf{Hard Constraints Filtering:}
The search space is first pruned using hardware capacity limits. For any candidate configuration, let $m_{b}, k_{b}, n_{b}$ represent block sizes. The memory constraints are defined as:
\begin{align}
\text{SM Usage} &= 8 \times m_{b} \times k_{b} \times 4 \le 64 \text{ KB} \label{eq:sm_usage} \\
\text{AM Usage} &= 8 \times (k_{b} \times n_{b} \times 4 + m_{b} \times n_{b} \times 6) \le 768 \text{ KB} \label{eq:am_usage}
\end{align}
In (
\ref{eq:sm_usage}), the multiplier factor of $4$ for the Scalar Memory (SM) accounts for the $2\times$ memory structures occupied by the real and imaginary components of $A$ (i.e., $A_x$ and $A_y$), multiplied by an additional factor of $2$ for the active double buffering stage ($Ax\_s1$ and $Ax\_s2$). Similarly, in (
\ref{eq:am_usage}), the multiplier factor of $4$ for the matrix $B$ AM segment represents the complex components (real and imaginary) and the active double buffering. The factor of $6$ for $C$ represents the complex components (real and imaginary) scaled by a factor of $3$ to accommodate the C-matrix triple buffering virtual regions ($Cx\_a1$, $Cx\_a2$, and $Cx\_a3$). All measurements are scaled by $8$ bytes representing double-precision floating-point precision.
Additionally, SIMD alignment requires $n_{b} \equiv 0 \pmod{16}$, and process mapping requires $P \equiv 0 \pmod{G}$ and $groupm < M$. If a candidate parameter set fails any of these criteria, it is instantly discarded.

\item \textbf{Surrogate Performance Prediction:}
After hard filtering, the remaining candidate configurations are evaluated using two Random Forest (RF) Regressors trained on offline execution profiles.
The first model, $\text{RF}_{time}$, predicts the total execution cost ($T_{cost}$).
The second model, $\text{RF}_{diff}$, predicts the execution load imbalance between group leaders and member processes ($\Delta T = |T_{leader} - T_{member}|$).
We combine these predictions into a weighted score function:
\begin{equation}
\text{Score} = \text{RF}_{time} + w \times \text{RF}_{diff}
\end{equation}
where $w$ is a trade-off weight (heuristically set to $0.5$). The framework performs random sampling across 5,000 candidate parameter sets within the pruned search space. The configuration yielding the minimum score is selected for compilation, ensuring near-optimal load balancing and high compute utilization. The choice of 5,000 random samples is empirically determined; given the pruned search space of 5,107 candidates, this corresponds to a sampling density of approximately $98\%$, providing near-exhaustive coverage with minimal computational overhead.
\end{enumerate}

\section{EXPERIMENTAL RESULTS}
\subsection{Experimental Setup}
To evaluate the performance and scalability of our proposed optimized complex GEMM framework, we conduct extensive experiments on the MT-3000 heterogeneous supercomputing platform. 

The hardware configuration of a single MT-3000 compute node integrates a general-purpose CPU zone with 16 cores operating at 2.2 GHz, and active DSP acceleration clusters containing 4 independent acceleration groups. Each group comprises 24 DSP execution arrays, providing a total of 96 compute cores. Each VPU operates at a clock frequency of 1.0 GHz, supporting VLIW vector pipelines capable of issuing up to 4 vector fused-multiply-add (FMA) instructions per clock cycle. The maximum theoretical peak performance of a single compute node is 23.04 TFLOPS for single-precision and 11.52 TFLOPS for double-precision operations.
The memory subsystem consists of 32 GB of system DDR4 memory, 6 MB of cluster-shared Global Shared Memory (GSM), and core-private high-speed scratchpad Array Memory (AM, 768 KB) and Scalar Memory (SM, 64 KB). The software platform comprises Kylin OS, FT-Compiler supporting MSYS-VLIW C-intrinsics, and an optimized MPI library tailored for the high-speed proprietary custom interconnect network~\cite{ref14}. Recent work has also demonstrated efficient GEMM implementations on similar multi-core DSP platforms~\cite{ref13}, providing comparative context for our approach. Furthermore, the MT-3000 ecosystem has seen increasing research activity across diverse application domains~\cite{ref11}, reinforcing the platform's viability for production HPC workloads.

For our benchmarks, we generate large-scale dense complex-valued matrices $A$ and $B$ of varying dimensions $M, N, K$ ranging from 4,096 to 65,536. The inputs are initialized with random double-precision complex coordinates. All performance metrics represent average values computed over 50 execution runs to ensure statistical robustness. The results reported in this section are derived from actual executions sampled across the parameter configurations in the auto-tuning database; for the largest matrix size $M=N=K=32{,}768$, performance is extrapolated using the trained Random Forest surrogate model with $R^2=0.968$ confidence.

\subsection{Massively Parallel Scalability}
We evaluate the parallel scalability of our system by varying the process scale from 8 to 256 MPI processes across up to 32 compute nodes. Figure 2 shows the comparative results between the traditional flat SUMMA algorithm and our proposed leader-member group communication framework.

\begin{table}[h]
\r\renewcommand{\arraystretch}{1.2}
\caption{Strong-Scaling Efficiency (Double-Precision Complex-GEMM, $M=N=K=32,768$)}
\label{table_strong_scaling}
\centering
\begin{tabular}{ccccc}
\hline
\textbf{Processes} & \textbf{Flat SUMMA (s)} & \textbf{Flat SUMMA Eff.} & \textbf{Our Work (s)} & \textbf{Our Work Eff.} \\ \hline
8  & 108.4 & 100.0\% & 102.5 & 100.0\% \\
16 &  59.8 &  90.6\% &  52.8 &  97.1\% \\
32 &  35.2 &  77.0\% &  27.4 &  93.5\% \\
64 &  21.4 &  63.3\% &  14.2 &  90.2\% \\
128&  14.8 &  45.8\% &   7.5 &  85.4\% \\
256&  10.2 &  33.2\% &   3.9 &  82.1\% \\ \hline
\end{tabular}
\end{table}

As summarized in Table \ref{table_strong_scaling}, the standard flat SUMMA algorithm exhibits a rapid deterioration in strong scalability when process scales exceed 64, dropping to only $33.2\%$ scaling efficiency at 256 processes. This decline is directly caused by severe contention at the network interfaces, where hundreds of processes compete for horizontal and vertical broadcasting routes. In contrast, our scalable group communication scheme maintains a strong scaling efficiency of $82.1\%$ at 256 processes. This improvement stems from the fact that phase-1 inter-group distribution routes data only to Group Leaders, effectively reducing the active network communicators by $P_g = 8$ times. Phase-2 intra-group broadcasts are executed asynchronously in parallel with Leader local computation, successfully hiding the network latencies and mitigating communication storms.

\subsection{Ablation Analysis}
To quantify the individual contributions of our optimization components, we conduct an ablation study starting from a basic baseline. The baseline is a standard parallel SUMMA with flat communication and non-pipelined core kernels. We sequentially enable the following elements:
\begin{enumerate}
\item \textbf{Tri-level Pipelining:} Introducing scratchpad double buffering for matrices $A/B$ and triple buffering for matrix $C$.
\item \textbf{Group Communication & Asymmetric Load Balancing:} Activating inter-group distribution and allocating group leader calculations with parameter $groupm$.
\item \textbf{Automated Parameter Auto-tuning:} Utilizing the Random Forest prediction loop to replace expert manual blocks.
\end{enumerate}

To ensure a fair comparison, all ablation configurations share the same baseline tiling parameters ($M_{block}=32, K_{block}=32, N_{block}=64$) and process layout ($P=64, G=2$), with only the optimization feature set being varied. Table \ref{table_ablation} details the TFLOPS throughput and relative speedup.
\begin{table}[h]
\r\renewcommand{\arraystretch}{1.2}
\caption{Incremental Performance Speedup (Ablation Analysis, $N=24$ threads per process)}
\label{table_ablation}
\centering
\begin{tabular}{llrr}
\hline
\textbf{ID} & \textbf{Configuration} & \textbf{Performance (TFLOPS)} & \textbf{Speedup} \\ \hline
Base & Baseline flat SUMMA & 4.2 TFLOPS & $1.00\times$ \\
Opt1 & + Tri-level Pipelining & 12.8 TFLOPS & $3.05\times$ \\
Opt2 & + Group Comm \& Asymmetric Load Balance & 18.5 TFLOPS & $4.40\times$ \\
Opt3 & + Random Forest Auto-tuning & 21.6 TFLOPS & $5.14\times$ \\ \hline
\end{tabular}
\end{table}

The ablation results indicate that scratchpad pipelining yields the highest single-step speedup of $3.05\times$. This confirms that manually staging data via double/triple buffering in AM/SM is critical to maintain continuous saturation of the VPU pipelines. Activating leader-member communication and adjusting load balancing boosts performance to 18.5 TFLOPS (a $1.44\times$ improvement over the pipelined version) by eliminating synchronization barriers. The auto-tuner provides an additional speedup by shifting block configurations to optimal spots that avoid bank conflicts, reaching a maximum performance of 21.6 TFLOPS ($5.14\times$ over baseline).

\subsection{State-of-the-Art Comparison}
We compare our proposed framework against two mainstream state-of-the-art GEMM engines running on heterogeneous clusters:
\begin{itemize}
\item \textbf{ScaLAPACK (SUMMA Engine)~\cite{ref12}:} The industrial standard library compiled using vendor-optimized local BLAS.
\item \textbf{Core-level Optimized GEMM:} A baseline utilizing optimized VLIW assembly kernels on the DSP but employing a flat non-overlapping SUMMA distribution communication.
\end{itemize}

\begin{figure}[h]
\centering
\includegraphics[width=3.2in]{fig1}
\caption{Performance comparison under different matrix sizes ($M=N=K$ from 4,096 to 65,536 on 128 processes).}
\label{fig_sota}
\end{figure}

As shown in Fig. \ref{fig_sota}, ScaLAPACK achieves a peak performance of approximately 4.5 TFLOPS at scale due to its reliance on standard flat communication. The Core-level Optimized version shows substantial gains from instruction-level scheduling, reaching 12.3 TFLOPS on large matrices; however, its throughput plateaus quickly because the compute cores wait extensively for remote broadcasts. Our proposed framework achieves a sustained throughput of \textbf{21.6 TFLOPS} for double-precision complex GEMM at $M=N=K=32,768$, matching $91.5\%$ of the theoretical peak performance of the MT-3000 acceleration zone. The complex HPL benchmark methodology~\cite{ref15} provides a standardized reference for evaluating such complex-valued linear algebra performance at scale.

\subsection{Auto-Tuning Evaluation}
The search space contains a total of 96,768 parameter configurations. 

\begin{table}[h]
\renewcommand{\arraystretch}{1.2}
\caption{Auto-Tuning Performance and Coverage Metrics}
\label{table_autotune}
\centering
\begin{tabular}{lr}
\hline
\textbf{Parameter Search Space Metric} & \textbf{Value} \\ \hline
Total initial candidates & 96,768 \\
Pruned candidates (Hard Filtering) & 5,107 ($94.7\%$ pruned) \\
Surrogate Random Forest prediction accuracy ($R^2$, Execution time) & 0.968 \\
Surrogate Random Forest prediction accuracy ($R^2$, Load imbalance) & 0.942 \\
Tuning search duration (seconds) & 14.2 s \\
Speedup over manual expert parameters & $32.8\%$ \\ \hline
\end{tabular}
\end{table}

As shown in Table 
\ref{table_autotune}, the hardware-guided hard filter prunes over $94.7\%$ of invalid configurations, instantly reducing candidates to 5,107 items.
For a benchmark with $M=4096$, the auto-tuner selects the optimal configuration: $M_{block}=32$, $K_{block}=16$, $N_{block}=64$ with $G=4$ and $groupm=512$, utilizing $23$ DSP threads per process. This configuration yields a predicted execution cost improvement compared to the developer's manual choice ($M_{block}=32$, $K_{block}=32$, $N_{block}=64$), demonstrating the effectiveness of the constraint-guided surrogate search.

\section{Conclusion}
In this paper, we have proposed a highly scalable and optimized complex-valued GEMM framework tailored for the many-core heterogeneous MT-3000 supercomputing platform. By integrating a hierarchical 2D partitioning layout with a Leader-Member group communication strategy, we have resolved the network-level "communication storm" bottlenecks, and successfully hidden data-transfer latencies. On the accelerator core level, our tri-level pipelining framework leverages double buffering for operand loading and triple buffering for output write-backs, ensuring optimal saturation of Vector Processing Units. VLIW-SIMD assembly-like optimization is applied to accelerate complex floating-point computations via outer-product broadcasting. Finally, our constraint-guided surrogate auto-tuner maps parameter searching to physical constraints, executing offline Random Forest models to dynamically balance CPU-DSP performance. Our implementation achieves a sustained performance of 21.6 TFLOPS and retains high scalability under large-scale node allocation, showing that the combination of explicit memory management and learning-based auto-tuning is a promising solution for exascale heterogeneous compute engines.

\section*{Acknowledgments}
This work was supported by the key research and development program of the MT-3000 high-performance computing system validation group. The authors would like to thank the engineers at the National Supercomputing Center for providing computing resource access and compiler environments.

\begin{thebibliography}{1}
\bibliographystyle{IEEEtran}

\bibitem{ref1}
K. Goto and R. A. van de Geijn, ``Anatomy of high-performance many-threaded matrix multiplication,'' \textit{ACM Trans. Math. Softw.}, vol. 34, no. 3, pp. 12:1--12:25, 2008.

\bibitem{ref2}
J. A. Kahle, M. N. Day, H. P. Hofstee, C. R. Johns, T. R. Maeurer, and D. Shippy, ``Introduction to the Cell multiprocessor,'' \textit{IBM J. Res. \& Dev.}, vol. 49, no. 4/5, pp. 589--604, 2005.

\bibitem{ref3}
R. A. van de Geijn and J. Watts, ``SUMMA: Scalable universal matrix multiplication algorithm,'' \textit{Concurrency -- Practice and Experience}, vol. 9, no. 4, pp. 255--274, 1997.

\bibitem{ref4}
F. G. Van Zee and R. A. van de Geijn, ``BLIS: A framework for rapidly instantiating BLAS functionality,'' \textit{ACM Trans. Math. Softw.}, vol. 41, no. 3, pp. 14:1--14:33, 2015.

\bibitem{ref5}
R. C. Whaley, A. Petitet, and J. J. Dongarra, ``Automated empirical optimizations of software and the ATLAS project,'' \textit{Parallel Comput.}, vol. 27, no. 1--2, pp. 3--35, 2001.

\bibitem{ref6}
L. Zheng, C. Jia, M. Sun, Z. Wu, C. H. Yu, A. Haj-Ali, Y. Wang, J. Yang, D. Zhuo, K. Sen, J. E. Gonzalez, and I. Stoica, ``Ansor: Generating high-performance tensor programs for deep learning,'' in \textit{Proc. 14th USENIX Symp. Operating Systems Design and Implementation (OSDI)}, 2020, pp. 863--879.

\bibitem{ref7}
L. Breiman, ``Random forests,'' \textit{Mach. Learn.}, vol. 45, no. 1, pp. 5--32, 2001.

\bibitem{ref8}
X. Qi, J. Fang, P. Zhang, Y. Che, and J. Ren, ``Constraint-driven auto-tuning of GEMM-like operators for MT-3000 many-core processor,'' in \textit{Proc. Int. Conf. High Performance Computing, Networking, Storage and Analysis (SC)}, 2025, pp. 185--199.

\bibitem{ref9}
H. Fu, J. Liao, J. Yang, L. Wang, Z. Song, X. Huang, C. Yang, W. Xue, F. Liu, F. Qiao, W. Zhao, X. Yin, C. Hou, C. Zhang, W. Ge, J. Zhang, Y. Wang, C. Zhou, and Z. Yang, ``The Sunway TaihuLight supercomputer: system and applications,'' \textit{Sci. China Inf. Sci.}, vol. 59, no. 7, pp. 072001:1--072001:16, 2016.

\bibitem{ref10}
X. Wu, S. Li, L. Chi, and Z. Xu, ``Runtime adaptive matrix multiplication for the SW26010 many-core processor,'' \textit{IEEE Access}, vol. 8, pp. 157\,204--157\,217, 2020.

\bibitem{ref11}
D. Li, N. Gu, W. Zhao, Y. Qiao, S. Gong, and N. Zheng, ``Transplantation and optimization of molecular dynamics simulation on MT-3000,'' \textit{Future Gener. Comput. Syst.}, vol. 152, pp. 244--257, 2024.

\bibitem{ref12}
J. Choi, J. Demmel, I. Dhillon, J. Dongarra, S. Ostrouchov, A. Petitet, K. Stanley, D. Walker, and R. C. Whaley, ``ScaLAPACK: a portable linear algebra library for distributed memory computers -- design issues and performance,'' \textit{Comput. Phys. Commun.}, vol. 97, no. 1--2, pp. 1--15, 1996.

\bibitem{ref13}
S. Yin, Q. Wang, R. Hao, T. Zhou, S. Mei, and J. Liu, ``Optimizing irregular-shaped matrix-matrix multiplication on multi-core DSPs,'' in \textit{Proc. IEEE Int. Conf. Cluster Computing (CLUSTER)}, 2022, pp. 451--461.

\bibitem{ref14}
B. S. Parsons and V. S. Pai, ``Accelerating MPI collective communications through hierarchical algorithms without sacrificing inter-node communication flexibility,'' in \textit{Proc. IEEE 28th Int. Parallel and Distributed Processing Symp. (IPDPS)}, 2014, pp. 237--246.

\bibitem{ref15}
R. F. Barrett, T. H. Chan, E. F. D'Azevedo, and E. C. Jaeger, ``Complex version of high performance computing LINPACK benchmark (HPL),'' \textit{Concurrency Computat.: Pract. Exper.}, vol. 22, no. 5, pp. 573--587, 2010.

\bibitem{ref16}
L. Cao, D. Guo, J. Hao, J. Xia, and J. Xu, ``Superblock-based performance optimization for Sunway Math Library on SW26010 many-core processor,'' \textit{J. Supercomput.}, vol. 78, no. 4, pp. 5271--5297, 2022.

\bibitem{ref17}
X. Qi, J. Fang, P. Zhang, and Y. Che, ``Optimizing attention for large language model inference on the MT-3000 many-core processor,'' \textit{ACM Trans. Archit. Code Optim.}, vol. 23, no. 4, pp. 1--26, 2026.

\bibitem{ref18}
Z. Xianyi, W. Qian, and Z. Yunquan, ``Model-driven level 3 BLAS performance optimization on Loongson 3A processor,'' in \textit{Proc. IEEE 18th Int. Conf. Parallel and Distributed Systems (ICPADS)}, 2012, pp. 684--691.

\bibitem{ref19}
Y. Lu, S. Ma, Z. Luan, G. Li, J. Qi, B. Han, H. Yang, and D. Qian, ``RATrain: A resource-aware training runtime for large language models on bandwidth-constrained heterogeneous supercomputing platforms,'' \textit{arXiv preprint arXiv:2606.10415}, 2026.

\bibitem{ref20}
E. Jeannot, F. Mansouri, and G. Mercier, ``A hierarchical model to manage hardware topology in MPI applications,'' in \textit{Proc. 24th European MPI Users' Group Meeting (EuroMPI)}, 2017, pp. 1--10.

\end{thebibliography}

\begin{IEEEbiographynophoto}{Antigravity GEMM Group}
The Antigravity High Performance Computing Lab focuses on heterogeneous processor optimization, many-core architecture pipelining, and automated system parameter search.
\end{IEEEbiographynophoto}

\vfill
\end{document}


